% 论文入口。沿用原 cumcmthesis.cls / cumcm2026.sty；原 example.tex 未覆盖。
% 在本目录运行 latexmk -xelatex main.tex 编译。
\documentclass[withoutpreface,bwprint]{cumcmthesis}
\usepackage{cumcm2026}
\usepackage{url}
\usepackage{placeins}
\hypersetup{hidelinks}
\setlength{\emergencystretch}{2em}
\renewcommand{\arraystretch}{1.13}
% 一级标题沿用中文数字；二、三级标题显示为4.1、4.2.1。
\renewcommand{\thesection}{\chinese{section}、}
\renewcommand{\thesubsection}{\arabic{section}.\arabic{subsection}}
\renewcommand{\thesubsubsection}{\thesubsection.\arabic{subsubsection}}
\title{非对称结算机制下的微网采购与储能协同调度}
\tihao{C}
\yearinput{2026}
\monthinput{9}
\dayinput{10}

\begin{document}
\maketitle
\begin{abstract}
针对微网采购、储能与预测更新的协同问题，建立以实际结算费用为中心的统一调度框架。由五倍紧急补购及1.5倍上调收费，推导无储能单时段的80\%日前采购分位数、70\%条件上调阈值及动态价格下的加权分位数，解释采购偏高的经济动机，再将该机制嵌入含功率、容量和效率限制的场景优化。

问题一建立购电与储能联合线性规划，并利用免费弃余条件证明存在不同时充放电的最优解。计算得到全天计划购电59482.70 kWh、购电费35126.95元，较无储能基准节省26.90\%，首末储电量均为6000 kWh。

问题二以一月实际账单校准有限候选预测器和执行器，冻结后对2月至12月334天顺序回测。主案保留此前28日完整残差轨迹，关闭额外风险与备用。总费用1401.28万元，较相同预测及初态下的点预测方案节省8.36\%，紧急购电85017.69 kWh。固定与自适应风险、备用单独及联合对照显示，附加保守机制并不普遍节费。

问题三将0、6、12、18时的合同修订与每10分钟的储能控制分开，在统一母模型外加入基于执行规则回放的购电候选筛选。各候选按相同历史场景、当前储电量、固定执行器及结算规则评分。四种筛选时点配置中，零时与日内均筛选的实际费用最低，为1387.15万元，较未筛选母模型降低0.747\%，但紧急购电由12647.76增至21029.90 kWh。原有状态、光伏及负荷信息来源对照仍按母模型报告。

问题四以同日配对的价格与净负荷场景回放候选，并按实际价格结算。日前版保留母模型，费用1479.41万元；滚动版采用全时点筛选，费用1467.02万元，较其母模型降低0.670\%，紧急电量为21293.41 kWh。日前两组筛选均未降低费用，故不替换原方案。预测器与执行器沿用冻结配置，筛选的适用问题及启用时点依据开发期回测确定，结果不构成独立测试或未来节费保证。正文主结果与附件5由对应最终执行数组统一生成。

\keywords{微网调度\quad 非对称结算\quad 场景优化\quad 执行回放\quad 候选筛选}
\end{abstract}

\input{sections/common_intro}
\FloatBarrier
\input{sections/question1}
\FloatBarrier
\input{sections/question2}
\FloatBarrier
\input{sections/question3}
\FloatBarrier
\input{sections/question4}
\FloatBarrier
\input{sections/model_evaluation}

\FloatBarrier
\CumcmAIUsed{建模方案整理、预测与调度程序生成、数据计算、文献检索核验、图表生成及论文初稿撰写}
\input{sections/references}

\clearpage
\begin{appendices}
\input{sections/05_appendix}
\end{appendices}
\end{document}
